\section{Problem Formulation}
\label{sec:formulation}

\subsection{Market Model and Data}
For a given stock, let \(t\) be the last observed session at the decision
cutoff. Indicator and chart tools use the 45 daily open, high, low, close and
volume (OHLCV) bars through \(t\). Alpha uses up to 600 observed bars through
the same cutoff. The 2022 Vietnam Securities Depository schedule~\cite{vsd2022settlement}
allows resale in the afternoon of T+2, counting T from the purchase date.
We use T+2.5 to denote the corresponding daily-bar convention
throughout the evaluation. LONG enters at the next observed Open
\(O_{t+1}\) and exits at the observed Close \(C_{t+3}\) two sessions after
entry. The session index counts observed bars. Calendar and execution
evidence are checked separately, and incomplete cycles have no label.

\subsection{Directional Prediction Task}
The target asks whether a candidate LONG would earn a positive net return.
With fee \(f=0.0025\) and slippage \(s=0.001\) applied on each leg, it is
\begin{equation}
  r_t = \frac{C_{t+3}}{O_{t+1}}
        \frac{(1-s)(1-f)}{(1+s)(1+f)}-1,
  \qquad
  y_t =
  \begin{cases}
    \mathrm{UP},   & r_t>0,\\
    \mathrm{DOWN}, & r_t\leq 0 .
  \end{cases}
  \label{eq:label}
\end{equation}
A LONG prediction maps to a buy followed by liquidation. SHORT maps to
CASH in the long-only ledger. For a complete account with non-overlapping
cycles and fractional execution at the reference prices, let \(a_t=1\)
for LONG and \(a_t=0\) for CASH. Its net return is
\(\prod_t(1+a_t r_t)-1\). Scenarios with lot sizes, delayed fills or pending
cash use their own ledger. The follow-up scores direction only where the
candidate LONG label is defined and compares account returns only where both
scheduled cycles meet the rule, even when a method chose CASH.

\subsection{Alpha Factor Formalisation and Adaptation}
\label{sec:alpha-representation}
The native registry combines 80 adaptations of
\emph{101 Formulaic Alphas}~\cite{kakushadze2016formulaic} with five
custom price--volume factors. A cross-sectional rank becomes a rank
against a stock's own trailing observations (20 bars by default).
Average daily volume uses recorded share volume, and the daily
typical-price proxy substitutes for intraday volume-weighted average price.
These changes alter what the source expressions measure. The registry
excludes 17 industry-neutralisation expressions, one market-capitalisation
expression and three price/volume expressions whose exclusion rationale is
undocumented. The custom formulas are given in \ref{app:signals}.

The native z-score--tanh transformation uses statistics from the supplied
prefix. Some handlers also fill earlier values or choose a volume proxy
from the supplied frame. The prefix ends at the decision cutoff, but its
later rows can still change an earlier diagnostic value. The causal
projection uses row-available operators and expanding normalization.

\subsection{Walk-Forward Alpha Selection}
The native rank uses the absolute Spearman information
coefficient (IC)~\cite{grinold1999active} between a factor and next-entry
net returns. A negative IC reverses the factor for the diagnostics and
report. Directional accuracy uses absolute scores above 0.05 and defaults
to 0.5 with fewer than ten such signals. The ranking Sharpe proxy uses
factor-signed net returns and population standard deviation. The account
ledger uses the actual LONG or CASH action. On the eligible candidate set \(\mathcal V\), let \(\mathcal N\)
be min--max scaling within a metric, with a value of 0.5 on a flat range.
The native score is
\begin{equation}
  S_k=0.40\,\mathcal N_{\mathcal V}(|\mathrm{IC}_k|)
      +0.35\,\mathcal N_{\mathcal V}(\mathrm{Acc}_k)
      +0.25\,\mathcal N_{\mathcal V}(\mathrm{Sharpe}_k).
  \label{eq:composite}
\end{equation}
Up to five candidates are reported after ranking the displayed scores.
Orientation and ranking reuse the same observations, so these are
selection diagnostics. The causal projection fits orientation on training
data and picks five factors and a threshold on purged validation data.
It freezes those choices and preprocessing before the test queries.

The earlier historical-context study paired the native Our system
projection with the common causal projection used by the BICL branches.
The follow-up uses the same causal alpha values for Our system and Common
and retains their distinct Decision prompts.

\subsection{Sentiment Signal Construction}
The Sentiment Agent uses the pinned
5CD-AI Vietnamese sentiment classifier~\cite{hf5cdsentiment}, built on
ViSoBERT~\cite{nguyen2023visobert}. For admitted articles \(i\), it
weights polarity \(\rho_i=P_i(\mathrm{POS})-P_i(\mathrm{NEG})\) by
publication age \(\Delta_i\) in days:
\begin{equation}
  S_t=\frac{\sum_i w_i\rho_i}{\sum_i w_i},
  \qquad w_i=2^{-\Delta_i/30}.
\end{equation}
The 90-day window requires three distinct, dated articles with accepted
scorer provenance and finite scores. Date-only publication is treated as
available the next day. The separate reliability flag requires at least
eight articles. Reported missingness means insufficient eligible evidence.
The historical provider supplies an aggregate and count, so peer-relative
fields on that path are placeholders. The follow-up supplies a missing-context record at each cutoff because
historical news availability could not be verified.

\subsection{Multi-Agent Decision Function}
QuantAgent~\cite{xiong2025quantagent} uses the shared Indicator, Pattern
and Trend reports. Our system adds the Alpha and Sentiment reports.
Component controls retain alpha only (A), sentiment only (S), or current
sentiment with past-donor alpha values (P). The backtest formats the
additional reports without an extra Alpha LLM call. The Decision Agent
receives the variant-specific prompt and emits LONG or SHORT. Historical
context branches can retain a completed parent action when their dependency
or revision rules apply. The scored prediction follows the final action,
and the ledger keeps the raw proposal separately.

BICL v5 estimates the probability of UP from separate historical and
target tasks. A Beta--Bernoulli teacher supplies probability draws
\(\theta\) for the utility check. With net-gain magnitude \(g\) and
loss magnitude \(\ell\),
\begin{equation}
  U(\mathrm{LONG})=\theta g-(1-\theta)\ell,\qquad
  U(\mathrm{CASH})=0,\qquad
  \Delta U=U(a^{\mathrm{proposed}})-U(a^0),
  \label{eq:revision_utility}
\end{equation}
where \(a^0\) is the Fact Anchor action. An eligible, fact-checked
revision replaces the anchor only if \(Q_{0.10}(\Delta U)>0\).
The utility quantile describes the specified teacher and magnitude model.

\subsection{Walk-Forward Evaluation Protocol}
Algorithm~\ref{alg:shared-reports} describes how the price-only and augmented
systems produce paired decisions and how BICL branches add context in the follow-up evaluation.

\begin{algorithm}[H]
\SetKwInOut{AlgInput}{Input}
\SetKwInOut{AlgOutput}{Output}
\caption{Walk-forward decisions with shared technical reports.}
\label{alg:shared-reports}
\AlgInput{Daily OHLCV, alpha registry, dated news, scheduled cutoffs,
          frozen historical resources and paired branch/repeat schedule.}
\AlgOutput{Recorded proposals and final actions, paired metrics and separate accounts.}
Reserve the evaluation cutoffs outside the historical bank\;
\ForEach{stock and decision cutoff \(t\) in chronological order}{
  Bound prices and dated news at \(t\); retrieve only cycles that ended before \(t\)\;
  Run Indicator, Pattern and Trend once for the shared upstream realization\;
  Prepare the profile-specific alpha and sentiment sections\;
  For follow-up retrieval, derive Trend and Pattern keys from cutoff OHLCV\;
  Freeze the reports and chart inputs; give each branch an independent copy\;
  \ForEach{Decision repeat and branch in dependency order}{
    Build the branch prompt from shared technical reports and its enabled evidence\;
    Add the eligible prior or task context for the corresponding BICL branch\;
    \eIf{the branch must reuse a parent action}{
      Record the same-repeat action and the reason for reuse\;
    }{
      Obtain a Decision proposal through retry and checkpoint handling\;
      Resolve the proposal under the branch policy and save the final action\;
    }
  }
}
After forecasts are locked, join eligible outcomes using Eq.~\ref{eq:label}\;
Score paired decisions and update separate LONG/CASH ledgers\;
\KwRet{the forecast records, account ledgers and paired metrics}\;
\end{algorithm}

The scheduled dates are evaluated chronologically with purging at the
train--validation and validation--test boundaries. Paired branches
reuse one upstream snapshot per point within each shared-report panel.
Historical context is frozen before the first query in a fold and excludes
all reserved evaluation points. Each retrieved cycle must have exited
strictly before the decision date. In the follow-up, forecasts are locked
before evaluation outcomes are computed. Section~\ref{sec:experimentalsetup}
gives the schedules, repetitions and account-level comparisons.
